\medskip

\begin{enumerate}
\item \textbf{Les droites (AB) et $(d)$ sont: sécantes non perpendiculaires.}

\medskip 

Un vecteur directeur de la droite (AB) est  le vecteur $\vect{\mathrm{AB}} \begin{pmatrix}  4\\4\\-2
\end{pmatrix}$.  Le vecteur $\vect{u} \begin{pmatrix}  2\\2\\-1
\end{pmatrix}$  est donc aussi un vecteur directeur de la droite (AB).

Un vecteur directeur de la droite $(d)$ est $\vect{u'} \begin{pmatrix}
3\\0\\- 5
\end{pmatrix}$.

Les vecteurs $\vect{u}$ et $\vect{u'}$ ne sont pas colinéaires donc les droites (AB) et $(d)$ ne sont pas parallèles.

Le repère est orthonormé, on peut donc utiliser les coordonnées des vecteurs pour calculer le produit scalaire.

 $\vect{u} \cdot \vect{u'} =2 \times 3 + 0 \times 2 - 1 \times (-5) = 11 \neq 0$

donc  les droites (AB) et $(d)$ ne sont pas perpendiculaires.

Une équation paramétrique de la droite (AB) de vecteur directeur $\vect{u}$ et passant par le point A$(-3;1;4)$ est :

$\begin{cases}x=-3+2 s\\ y=1 +2s \\ z=4-s\end{cases} \quad \text{ où } s \in \R$.

Les deux droites sont sécantes si et seulement si le système suivant admet une unique solution :

$\begin{aligned} \begin{cases}- 3 + 2 s=-6+3t\\ 1 + 2s=1 \\ 4-s =9-5t\end{cases} & \iff \begin{cases}2 s -3t =-3\\ s=0 \\ s-5t=-5\end{cases} &
& \iff \begin{cases}-3t =-3\\ s=0 \\ -5t=-5\end{cases}
& \iff\begin{cases}t =1\\ s=0 \\ t=1\end{cases}\\
\end{aligned}$

Le système admet une unique solution donc les droites sont sécantes. 

\item \textbf{La droite(AB) est orthogonale au plan $\mathcal{P}$.}

Un vecteur normal au plan $\mathcal{P}$ est $\vect{n}\begin{pmatrix}4\\4\\-2\end{pmatrix}$ soit le vecteur $\vect{AB}$.

La droite (AB) est donc orthogonal au plan $\mathcal{P}$.

\item \textbf{Les plans $\mathcal{P}$ et $\mathcal{P}'$ sont perpendiculaires.}

Un vecteur normal au plan $\mathcal{P'}$ est $\vect{n'}\begin{pmatrix}
2\\1\\6
\end{pmatrix}$.

Les vecteurs $\vect{n}$ et $\vect{n'}$ ne sont pas colinéaires donc les plans ne sont pas parallèles.

On est dans un repère orthonormé, on peut donc utiliser les coordonnées pour calculer le produit scalaire.

$\vect{n} \cdot \vect{n'} = 4  \times 2+ 4\times 1-2 \times (-6) = 8+4-12=0$

Les vecteurs normaux sont orthogonaux donc les plans sont perpendiculaires.
\item  \textbf{La valeur de l'angle $\widehat{\mathrm{B A C}}$ arrondie au degré est  $51\degres$.}

On a : $\vect{\mathrm{AC}} \begin{pmatrix}
	3 \\ 0 \\ -5
\end{pmatrix}$ et $\vect{\mathrm{AB}} \begin{pmatrix}  4\\4\\-2
\end{pmatrix}$.

On est dans un repère orthonormé donc on peut utiliser les coordonnées pour calculer les longueurs et le produit scalaire.

Calculons le produit scalaire de ces deux vecteurs de deux manières différentes :

d'une part : $\vect{\mathrm{AC}}\cdot \vect{\mathrm{AB}} = 3\times4 + 0 \times 4 -5 \times (-2)= 22$

d'autre part :

AC$=\sqrt{(-3)^2+0^2+ (-5)^2}=\sqrt{34}$ et AB$=\sqrt{4^2+4^2+(-2)^2} =\sqrt{36} = 6$

On a donc  $\vect{\mathrm{AC}}\cdot \vect{\mathrm{AB} }= \text{AB} \times \text{AC} \times \cos \left(\widehat{\text{BAC}} \right)= 6\sqrt{34} \cos \left(\widehat{\text{BAC}} \right)$

D'où : $ 22 = 6\sqrt{34} \cos \left(\widehat{\text{BAC}}\right)$, 
donc $ \cos \left(\widehat{\text{BAC}} \right) =\dfrac{22}{6\sqrt{34}} $

D'où $ \widehat{\text{BAC}} \approx 51\degres$
\end{enumerate}
\medskip
